% ACMSE 2027 research poster — DOUBLE-ANONYMOUS REVIEW DRAFT.
% IMPORTANT: Before submission, compile using ACMSE 2027's modified ACM Master
% Article Template, starting from its sample-sigconf.tex (not stock acmart).
% Template: https://acmse.net/2027/papers-format/
% Data source: AivoBlock v0.1.0-research-prototype; CI run 39740779280,
% commit ade6ff42737d972658e9c309f9a2d8cffbfb988c.
% Do not add author, institution, repository identity, acknowledgments or
% identifying metadata to the review PDF without venue-chair guidance.
\documentclass[sigconf,anonymous]{acmart}

% Review-only metadata. Final publication identifiers come from ACM if accepted.
\setcopyright{none}
\acmConference[ACMSE '27]{2027 ACM Southeast Conference}{April 15--17, 2027}{Decatur, Alabama, USA}
\acmYear{2027}
\copyrightyear{2027}
% No fabricated DOI/ISBN in anonymous review metadata; use ACM-issued values after acceptance.
\hypersetup{pdfauthor={},pdfsubject={Anonymous ACMSE research poster abstract}}

\begin{document}
\title{AivoBlock: Reproducible Verification of Partitioned AI Inference in a Permissioned Blockchain}

\begin{abstract}
We investigate whether small artificial-intelligence inference tasks can be executed by distributed workers, checked deterministically by blockchain validators, finalized, and rewarded within a permissioned network. AivoBlock combines Python workers performing row-partitioned integer neural-network inference, Rust validators that recompute submitted results, and deterministic Haskell ledger transitions. A fixed-validator quorum finalizes the resulting blocks. In an archived experiment, 90 trained-model task requests across three fresh three-validator testnets had a pooled median client-observed request-to-finality duration of 730.760 ms with 100 ms status polling. A separate single-process microbenchmark measured 0.742 ms for direct integer inference and 3.412 ms for Rust full-block verification on a 256-row, 64-partition workload. These measurements have different boundaries and must not be combined into an end-to-end overhead estimate. The prototype shows integration and reproducible characterization for constrained workloads; it does not establish permissionless proof of useful work or efficient verification of arbitrary AI computation.
\end{abstract}

\ccsdesc[500]{Computer systems organization~Distributed architectures}
\ccsdesc[300]{Computing methodologies~Neural networks}
\keywords{Permissioned Blockchain, Deterministic Inference, Verifiable Computation, Distributed AI, Reproducibility}
\maketitle

\section{Motivation and Research Question}
Replacing blockchain hash puzzles with useful computation has been explored in prior work, including neural-network training proposals \cite{baldominos2019coinai,lihu2020useful} and formally motivated useful-work protocols \cite{fitzi2022ofelimos}. A recent inference-based proposal also analyzes permissionless mining security \cite{wong2026poml}. Our work does \emph{not} claim a new permissionless proof-of-work mechanism. Instead, we ask a narrower systems question: can a fixed validator committee agree on results and reward settlement for small, partitioned AI \emph{inference} tasks while avoiding nondeterministic floating-point outputs? The contribution is a concrete cross-language integration and an archived, scope-limited experimental characterization.

\section{Prototype and Experimental Method}
A signed network configuration fixes validator and worker membership. An authorized client commits to an integer model, input rows, and work partitions. Python workers return signed result claims. Rust validators authenticate those claims, check commitments, and \emph{fully recompute} the integer inference rather than verifying a succinct proof. Validators cross-check deterministic reward settlement against the Haskell core and finalize through a voting-power quorum above two thirds. The fixed committee provides no public Sybil resistance or task-based mining difficulty. This separation of task correctness from ledger agreement is deliberate: returning a valid inference output alone does not justify a block or imply permissionless consensus security.

A floating-point model is trained off-chain on synthetic glyph data and quantized into a specified QMLP-V1 integer representation. Exact execution, not agreement between floating-point kernels, defines eligible consensus results. The released evidence contains 270 client-observed chain tasks across three 90-task conditions (each condition: three fresh three-validator testnets with 30 tasks each). The trained-model condition uses 12 input rows per task and 100 ms client status polling. Separately, three invocations of a single-process 256-row/64-partition QMLP microbenchmark compare direct computation with verification. All results are descriptive and tied to one archived GitHub Actions runner environment; the separate-runner audit checks archive integrity and reanalyzes recorded statistics but does not independently repeat the live network experiment. The evidence identifies model, task and finalized-block commitments, checks unique request identifiers and finalized heights, and archives timing records for reanalysis. Pooled observations within a shared testnet launch are dependent; we report descriptive statistics rather than treating the 90 tasks as independent network deployments.

\section{Observed Results}
On the synthetic base-model test set, the quantized inference achieved 94.8\% accuracy; tested integer-ONNX/reference predictions had zero mismatches in 500 examples. In the trained-model chain condition, all 90 archived task records were finalized and the pooled median \emph{client-observed} request-to-finality duration was 730.760 ms (nearest-rank p95: 779.342 ms). These timings include request submission, protocol activity, and polling, and are not consensus-only latency. For the distinct 256-row/64-partition microbenchmark, direct integer inference averaged 0.742 ms, Rust full-block verification averaged 3.412 ms, and persistent Haskell cross-checked verification averaged 8.282 ms. The latter paths perform additional checks; none of these microbenchmark values is directly comparable to the observed chain finality duration. The CPU-only service measured a 1.880 ms median over 90 separately timed executions, but this excludes peer communication and consensus, so no service-to-chain overhead ratio is claimed.

\section{Verification Cost and Interpretation}
On a fixed 256-row inference fixture, increasing the number of partitions from one to 64 increased Rust verification time from approximately 0.566 ms to 3.404 ms, reflecting additional claim/signature and bookkeeping costs. Persistent Haskell settlement cross-checking averaged 8.282 ms in the 64-partition benchmark, while starting a fresh Haskell process for each call averaged 18.757 ms. These values describe single-process verification choices, not performance scaling from distributing AI work. Full validator recomputation remains the dominant design limitation: partitioning work across registered workers does not make checking arbitrary results asymptotically cheaper. The measurements motivate experiments with batched claims or selectively checked computation, while retaining an explicit adversarial correctness model.

\section{Relationship to Prior Work}
\emph{Coin.AI} and earlier useful-work designs investigate attaching machine-learning tasks to mining incentives; \emph{Ofelimos} instead provides a security-analyzed optimization-based consensus design \cite{baldominos2019coinai,fitzi2022ofelimos}. Recent proof-of-ML-inference research examines a permissionless security construction \cite{wong2026poml}. AivoBlock differs primarily in its implemented, small-scale, permissioned integer-inference pipeline and released descriptive observations. It does not establish originality of inference-based mining, a new Byzantine agreement theorem, or lower-energy consensus. This distinction limits the claim to an evaluated engineering artifact rather than a novel general-purpose consensus primitive.

\section{Limitations and Discussion}
Validator recomputation imposes substantial work and does not deliver succinct proof-of-computation benefits. Small synthetic datasets, a fixed authority, three- or four-validator testbeds, and one CI hardware environment prevent claims about geographic replication, open membership, production scalability, energy savings, or incentive security. The next experiment should compare against a non-blockchain distributed service with aligned trust and timing boundaries, separate network/consensus instrumentation, and independent hardware reproduction. The present results support a reproducible \emph{permissioned prototype}, not a solution to permissionless useful-work consensus.

\bibliographystyle{ACM-Reference-Format}
\bibliography{references}
\end{document}
